\exercisesection{4}

\begin{enumerate}
\def\labelenumi{\arabic{enumi}.}
\item
  Let \((\mathbf{X}_{\alpha}, f_{\alpha\beta})\) be an inverse system of topological spaces, whose indexing set is directed, \(X = \varprojlim X_{\alpha}\) and \(f_{\alpha}\) be the canonical mapping from X to \(X_{\alpha}\) . Suppose that, for all \(\alpha, f_{\alpha}\) is surjective; show that, if the \(X_{\alpha}\) are irreducible, X is irreducible (cf.~General Topology, Chapter I, §4, no. 4, Corollary to Proposition 9). In particular, every product of irreducible spaces is irreducible.
\item
  A point \(x\) in an irreducible space \(\mathbf{X}\) is called generic if \(\{x\}\) is everywhere dense in \(\mathbf{X}\) .
\end{enumerate}

\begin{enumerate}
\def\labelenumi{(\alph{enumi})}
\tightlist
\item
  If \(\mathbf{X}\) is a Kolmogoroff space (General Topology, Chapter I, § 1, Exercise 2), it admits at least one generic point; if \(\mathbf{X}\) is an accessible space (General
\end{enumerate}

Topology, Chapter I, § 8, Exercise 1), it cannot admit a generic point unless it consists of a single point.

\begin{enumerate}
\def\labelenumi{(\alph{enumi})}
\setcounter{enumi}{1}
\item
  Give an example of an infinite irreducible accessible space (cf.~General Topology, Chapter I, § 8, Exercise 5).
\item
  Let X, Y be two irreducible spaces each admitting a generic point; suppose also that Y admits only one generic point y. Let \(f: X \to Y\) be a continuous mapping; for \(f(X)\) to be dense in Y, it is necessary and sufficient that, for every generic point x of X, \(f(x) = y\) .
\item
  Let \((\mathbf{X}_{\alpha}, f_{\alpha \beta})\) be an inverse system of irreducible spaces whose indexing set is directed; suppose that each of the \(\mathbf{X}_{\alpha}\) admits a single generic point \(x_{\alpha}\) . Show that, if, for \(\alpha \leqslant \beta, f_{\alpha \beta}(\mathbf{X}_{\beta})\) is dense in \(\mathbf{X}_{\alpha}\) , then \(\mathbf{X} = \varprojlim \mathbf{X}_{\alpha}\) is irreducible and admits a single generic point (same method as in Exercise 1).
\end{enumerate}

\begin{enumerate}
\def\labelenumi{\arabic{enumi}.}
\setcounter{enumi}{2}
\item
  Let X be a topological space and \((Y_{\alpha})\) a right directed family of subspaces of X; show that if each of the subspaces \(Y_{\alpha}\) is irreducible and X is the union of the family \((Y_{\alpha})\) , X is irreducible. Deduce an example where X is a Kolmogoroff space, each of the \(Y_{\alpha}\) admits a generic point, but X does not admit a generic point.
\item
  Let Y be an irreducible space admitting a single generic point y, X a topological space and \(f: X \rightarrow Y\) a continuous mapping.
\end{enumerate}

\begin{enumerate}
\def\labelenumi{(\alph{enumi})}
\item
  For every irreducible component \(\mathbf{Z}\) of \(\mathbf{X}\) meeting \(f^{-1}(y), f(\mathbf{Z})\) is dense in \(\mathbf{Y}\) .
\item
  Give an example where \(f(\mathbf{X})\) is dense in \(\mathbf{Y}\) but \(f^{-1}(y)\) is empty (take \(\mathbf{X}\) to be a subspace of \(\mathbf{Y}\) ).
\item
  Show that, if \(\mathbf{Z}\) admits a generic point \(z\) and \(f(\mathbf{Z})\) is dense in \(\mathbf{Y}\) , then \(f(z) = y\) and \(z\) is a generic point of \(\mathbf{Z} \cap f^{-1}(y)\) .
\end{enumerate}

\begin{enumerate}
\def\labelenumi{\arabic{enumi}.}
\setcounter{enumi}{4}
\item
  Let G be a connected semi-topological group (General Topology, Chapter III, § 1, Exercise 2). Show that, if G only admits a finite number of irreducible components, it is necessarily irreducible (observe that the irreducible components are derived from the irreducible component containing the identity element e by left or right translation).
\item
  Let \(\mathbf{X}\) be a topological space.
\end{enumerate}

\begin{enumerate}
\def\labelenumi{(\alph{enumi})}
\item
  Let x be a point of X and U an open neighbourhood of x with only a finite number of irreducible components. Show that there exists a neighbourhood V of x such that every neighbourhood contained in V is connected.
\item
  Suppose that every point of X possesses an open neighbourhood with only a finite number of irreducible components. Show that the following properties are equivalent:
\end{enumerate}

(α) The irreducible components of X are open.

(β) The irreducible components of X are identical with its connected components.

(γ) The connected components of X are irreducible.

(δ) Two distinct irreducible components of X do not meet.

\begin{enumerate}
\def\labelenumi{\arabic{enumi}.}
\setcounter{enumi}{6}
\item
  Let X be a compact metrizable space and R an open equivalence relation on X such that the quotient space X/R is irreducible. Show that there exists a point \(x \in X\) whose class mod. R is everywhere dense. (Note first that every open subset of X which is saturated with respect to R is everywhere dense. Then use Baire's Theorem.)
\item
  Show that the product of two Noetherian spaces is Noetherian. (If X and Y are Noetherian spaces and A is an open subset of \(X \times Y\) , show that, for all \(x \in pr_{1}A\) , there exists an open neighbourhood V of x such that \(V \times A(x) \subset A\) and deduce that A is quasi-compact.)
\item
  \begin{enumerate}
  \def\labelenumii{(\alph{enumii})}
  \tightlist
  \item
    Show that the prime spectrum of a commutative ring is a Kolmogoroff space (General Topology, Chapter I, § 1, Exercise 2), in which every irreducible component admits a (unique) generic point.
  \end{enumerate}
\end{enumerate}

\begin{enumerate}
\def\labelenumi{(\alph{enumi})}
\setcounter{enumi}{1}
\tightlist
\item
  Let A be an integral domain and \(\xi = (0)\) the generic point of \(\mathbf{X} = \operatorname{Spec}(\mathbf{A})\) . For the point \(\xi\) to be isolated in X, it is necessary and sufficient that the intersection of the prime ideals \(\neq 0\) of A be an ideal \(\neq 0\) . Give an example of a local ring with this property.
\end{enumerate}

\begin{enumerate}
\def\labelenumi{\arabic{enumi}.}
\setcounter{enumi}{9}
\item
  Let A be a ring and N its nilradical. For the spectrum of A to be discrete, it is necessary and sufficient that A/N be the direct composition of a finite number of fields.
\item
  Let K be a field, A the polynomial ring K{[}X, Y{]} and B the quotient ring A/(XY); show that Spec(B) is connected but has two distinct irreducible components.
\item
  Let Y be a completely regular space, \(\mathbf{A} = \mathcal{C}^{\infty}(\mathbf{Y}; \mathbf{R})\) and \(\mathbf{B} = \mathcal{C}(\mathbf{Y}; \mathbf{R})\) . Show that the Stone-Čech compactification \(\beta Y\) of Y is canonically identified with an everywhere dense subspace of \(\operatorname{Spec}(A)\) and with an everywhere dense subspace of \(\operatorname{Spec}(B)\) .
\item
  Let \(\mathbf{A}\) be a ring and \(\mathbf{X} = \operatorname{Spec}(\mathbf{A})\) its spectrum.
\end{enumerate}

\begin{enumerate}
\def\labelenumi{(\alph{enumi})}
\item
  If a subset F of X is closed, it has the following two properties: (α) for all p ∈ F, V(p) ⊂ F; (β) for all p ∉ F, there exists a closed subset of V(p) which contains F ∩ V(p) and does not contain p.
\item
  Suppose X is Noetherian. Show that every subset F of X with properties \((\alpha)\) and \((\beta)\) of (a) is closed in X (consider the irreducible components of \(\overline{F}\) and use Exercise 9 (a)).
\item
  With the notation of Exercise 12, take Y to be the interval {[}0, 1{]} of R. Show that Y is not closed in X = Spec(A) (cf.~§ 2, Exercise 15) but satisfies conditions ( \(\alpha\) ) and ( \(\beta\) ) of (a).
\end{enumerate}

\begin{enumerate}
\def\labelenumi{\arabic{enumi}.}
\setcounter{enumi}{13}
\tightlist
\item
  Let A be a ring, \(\mathbf{X} = \operatorname{Spec}(\mathbf{A})\) its spectrum and P the set of minimal prime ideals of A.
\end{enumerate}

\begin{enumerate}
\def\labelenumi{(\alph{enumi})}
\item
  Let \(R_{\mathfrak{p},\mathfrak{p}^{\prime}}\) be the symmetric and reflexive relation: ``there exists a prime ideal \(\mathfrak{p}''\) of A containing \(\mathfrak{p} + \mathfrak{p}'''\) between the elements \(\mathfrak{p},\mathfrak{p}'\) of P; let S be the equivalence relation whose graph is the smallest of the graphs of equivalences containing the graph of R (Set Theory, Chapter II, § 6, Exercise 10). Show that, if I is an equivalence class with respect to S, the set \(V_{I} = \bigcup_{\mathfrak{p}\in I}V(\mathfrak{p})\) is connected.
\item
  Show that, if P is finite (and in particular if A is Noetherian), the \(V_{I}\) are the connected components of X. Does this result extend to the case where P is infinite (cf.~Exercise 12)?
\end{enumerate}

¶ 15. Let A be a ring and a a finitely generated ideal of A. Show that the following properties are equivalent: (α) \(a^{2} = a\) ; (β) A is generated by an idempotent; (γ) V(a) is both open and closed in X = Spec(A) and a is minimal among the ideals b such that r(b) = r(a) (use § 2, Corollary 3 to Proposition 4). Give an example where \(a^{2} = a\) , a is not finitely generated and does not contain an idempotent (cf.~§ 2, Exercise 15).

¶ 16. (a) Let A be an absolutely flat (commutative) ring (Chapter I, §2, Exercise 17). Show that X = Spec(A) is a totally disconnected compact space, that every prime ideal p of A is maximal and that \(A_{p}\) is canonically isomorphic to the field A/p (show first that, for all f ∈ A, V(f) is both open and closed in X; on the other hand use §3, Exercise 9).

\begin{enumerate}
\def\labelenumi{(\alph{enumi})}
\setcounter{enumi}{1}
\tightlist
\item
  Show that the mapping \(\mathfrak{a}\mapsto\mathrm{V}(\mathfrak{a})\) is a bijection of the set of ideals of A onto the set of closed subsets of X. Moreover, the following conditions are equivalent: ( \(\alpha\) ) V(a) is open; ( \(\beta\) ) a is finitely generated; ( \(\gamma\) ) a is generated by an idempotent; ( \(\delta\) ) A/a is a projective A-module.
\end{enumerate}

* (c) Let \(\tilde{A}\) be the structure sheaf of rings over \(X\) . Show that every \(\tilde{A}\) -Module \(\mathcal{F}\) is of the form \(\tilde{E}\) , where \(E\) is an A-module defined up to a unique isomorphism (take \(E = \Gamma(X, \mathcal{F})\) and use the fact that every point of \(X\) admits a fundamental system of neighbourhoods which are both open and closed). For every ideal \(a\) of \(A\) , show that \(\operatorname{Hom}_A(a, E)\) is identified with \(\Gamma(X - V(a), \tilde{E})\) (note that \(X - V(a)\) is the union of the \(X_e\) , where \(e\) runs through the set of idempotents \(e \in a\) ). Deduce that, for \(E\) to be an injective A-module, it is necessary and sufficient that \(\tilde{E}\) be flabby. *

\begin{enumerate}
\def\labelenumi{(\alph{enumi})}
\setcounter{enumi}{3}
\tightlist
\item
  Let A be a ring and \(\mathfrak{N}\) its nilradical. The following properties are equivalent:
\end{enumerate}

(α) \(A/\mathfrak{N}\) is absolutely flat.

(β) \(\mathbf{X} = \operatorname{Spec}(\mathbf{A})\) is Hausdorff.

(γ) Every point of X is closed (in other words, every prime ideal of A is maximal).

(Use Exercises 9 of § 3.)

¶ 17. * (a) Let X be a totally disconnected compact space and O a sheaf of rings over X such that, for all \(x \in X\) , \(\mathcal{O}_x\) is a commutative field. Show that, if A = \(\Gamma(X, \mathcal{O})\) , the ring A is absolutely flat and its spectrum (with the sheaf of rings \(\tilde{A}\) ) is canonically identified with X (with the sheaf O). (Observe that for all \(f \in A\) the set of \(x \in X\) such that \(f_x = 0\) is both open and closed in X and deduce the existence of \(g \in A\) such that \(f = gf^2\) .) *

\begin{enumerate}
\def\labelenumi{(\alph{enumi})}
\setcounter{enumi}{1}
\item
  Let X be a totally disconnected compact space, K a (commutative) field and A the ring of locally constant mappings from X to K. Show that A is absolutely flat (argue directly or use (a)).
\item
  Let A be the ring of real-valued step functions defined on I = {[}0, 1{[}, whose points of discontinuity are of the form \(k/2^n\) ( \(n \in \mathbf{N}, 0 \leqslant k < 2^n\) ) and which are right continuous. Let \(\mathfrak{S}\) be the set of finite unions of half-open intervals \([x, y[ \subset I\) whose extremities are of the form \(k/2^n\) . For every \(f \in A\) ,
\end{enumerate}

the set \(Z(f) = f^{-1}(0)\) belongs to \(\mathfrak{S}\). If a is an ideal of A, the set \(\mathfrak{B}\) of \(Z(f)\) , where f runs through a, is such that: (1) the intersection of two sets of \(\mathfrak{B}\) belongs to \(\mathfrak{B}\); (2) every set of \(\mathfrak{S}\) containing a set of \(\mathfrak{B}\) belongs to \(\mathfrak{B}\). Obtain the converse. For a to be prime, it is necessary and sufficient that \(\mathfrak{B}\) be the base of a filter which converges to a point of {[}0, 1{]}. Show that A is an absolutely flat reduced ring and that, for every maximal ideal m of A, \(A_{m}\) is isomorphic to the field R. Describe the space Spec(A).

¶ 18. (a) Let Y be a completely regular space. Show that the following conditions are equivalent:

\((\alpha)\) Every prime ideal of \(\mathcal{C}(\mathbf{Y};\mathbf{R})\) is maximal.

(β) Every countable intersection of open subsets of Y is open.

(γ) Every continuous real-valued function on Y is locally constant.

(δ) The ring \(\mathcal{C}(\mathbf{Y};\mathbf{R})\) is absolutely flat.

(Use Exercise 15 (f) of § 2.)

If these hold, the space \(\mathbf{X} = \operatorname{Spec}(\mathcal{C}(\mathbf{Y};\mathbf{R}))\) is canonically identified with the Stone-Čech compactification \(\beta \mathbf{Y}\) of \(\mathbf{Y}\) . Such a space \(\mathbf{Y}\) is called flat.

\begin{enumerate}
\def\labelenumi{(\alph{enumi})}
\setcounter{enumi}{1}
\item
  Every subspace of a flat space is flat. Every completely regular quotient space of a flat space is flat. Every finite product of flat spaces is flat. Every sum of flat spaces is a flat space.
\item
  In a flat space, every countable subspace is closed and discrete. In particular every countable flat space is discrete.
\item
  Every Weierstrassian completely regular space (General Topology, Chapter IX, § 1, Exercise 22) (and a fortiori every compact space), which is flat, is finite.
\item
  The space associated with a non-trivial ultrafilter is extremely disconnected (General Topology, Chapter I, § 11, Exercise 21) but is not flat.
\item
  The discrete space \(\mathbf{N}\) is flat but its Stone-Čech compactification is not flat and the spectra of the rings \(\mathcal{C}(\mathbf{N};\mathbf{R})\) and \(\mathcal{C}^{\infty}(\mathbf{N};\mathbf{R})\) are therefore not isomorphic.
\end{enumerate}

\begin{enumerate}
\def\labelenumi{\arabic{enumi}.}
\setcounter{enumi}{18}
\tightlist
\item
  \begin{enumerate}
  \def\labelenumii{(\alph{enumii})}
  \tightlist
  \item
    Let \(\phi: A \to B\) be a ring homomorphism. Show that, for every ideal \(\mathfrak{b}\) of \(B\) , \(\overline{^a\phi(V(b))} = V(\overline{\phi}^{-1}(b))\) .
  \end{enumerate}
\end{enumerate}

\begin{enumerate}
\def\labelenumi{(\alph{enumi})}
\setcounter{enumi}{1}
\item
  Deduce from (a) that, for \(^{a}\phi(\text{Spec}(B))\) to be dense in \(\text{Spec}(A)\) , it is necessary and sufficient that \(\text{Ker}(\phi)\) be a nil ideal of A.
\item
  Show that, if every \(b \in \mathbf{B}\) may be written as \(b = h\phi(a)\) , where \(h\) is invertible in \(\mathbf{B}\) and \(a \in \mathbf{A}\) , \(^a\phi\) is a homeomorphism of \(\operatorname{Spec}(\mathbf{B})\) onto a subspace of \(\operatorname{Spec}(\mathbf{A})\) .
\end{enumerate}

\begin{enumerate}
\def\labelenumi{\arabic{enumi}.}
\setcounter{enumi}{19}
\item
  Give an example of a ring homomorphism \(\phi: A \to B\) such that every maximal ideal of \(A\) is of the form \(\bar{\phi}(m)\) , where \(m\) is a maximal ideal of \(B\) , but there exist prime ideals of \(A\) which are not inverse images under \(\phi\) of ideals of \(B\) (take \(B\) to be a field).
\item
  Let \((\mathbf{A}_{\alpha}, \phi_{\beta \alpha})\) be a direct system of rings whose indexing set is directed, \(\mathbf{A} = \varinjlim \mathbf{A}_{\alpha}\) and \(\phi_{\alpha}: \mathbf{A}_{\alpha} \to \mathbf{A}\) the canonical homomorphism. If we write \(\mathbf{X}_{\alpha} = \operatorname{Spec}(\mathbf{A}_{\alpha})\) , \((\mathbf{X}_{\alpha}, {}^{a}\phi_{\beta \alpha})\) is an inverse system of topological spaces and, if \(\mathbf{X} = \operatorname{Spec}(\mathbf{A})\) , the \({}^{a}\phi_{\alpha}: \mathbf{X} \to \mathbf{X}_{\alpha}\) form an inverse system of continuous mappings. Show that \(u = \varprojlim {}^{a}\phi_{\alpha}\) is a homeomorphism of \(\mathbf{X}\) onto \(\varprojlim {}^{\infty}\mathbf{X}_{\alpha}\) . (Prove
\end{enumerate}

first that, if, for all \(\alpha\) , \(p_{\alpha}\) is a prime ideal of \(\mathbf{A}_{\alpha}\) such that \(p_{\alpha} = \phi_{\beta \alpha}^{-1}(p_{\beta})\) for all \(\alpha \leqslant \beta\) , the union \(p\) of the \(\phi_{\alpha}(p_{\alpha})\) is a prime ideal of \(\mathbf{A}\) ; conversely, every prime ideal \(p\) of \(\mathbf{A}\) is the union of the \(\phi_{\alpha}(p_{\alpha})\) , where \(p_{\alpha} = \overline{\phi}_{\alpha}(p)\) . Finally, if \(f_{\alpha} \in \mathbf{A}_{\alpha}\) and \(f = \phi_{\alpha}(f_{\alpha})\) , note that the relation \(p \in \mathbf{X}_f\) is equivalent to

\[
u (\mathfrak {p}) \in \overline {{\operatorname{pr}}} _ {\alpha} ^ {- 1} ((\mathrm{X} _ {\alpha}) _ {f _ {\alpha}}).
\]

\begin{enumerate}
\def\labelenumi{\arabic{enumi}.}
\setcounter{enumi}{21}
\tightlist
\item
  \begin{enumerate}
  \def\labelenumii{(\alph{enumii})}
  \tightlist
  \item
    Let M be an A-module and \((\mathbf{N}_{i})\) a finite family of submodules of M. Show that \(\operatorname{Supp}\left(\mathbf{M}/\bigcap_{i}\mathbf{N}_{i}\right)=\bigcup_{i}\operatorname{Supp}(\mathbf{M}/\mathbf{N}_{i})\) (note that \(M/\bigcap_{i}N_{i}\) is isomorphic to a submodule of the direct sum of the \(M/N_{i}\) ).
  \end{enumerate}
\end{enumerate}

\begin{enumerate}
\def\labelenumi{(\alph{enumi})}
\setcounter{enumi}{1}
\item
  Let \(p\) be a prime number; the support of the \(\mathbf{Z}\) -module \(\mathbf{Z}\) is distinct from the closure of the union of the supports of the \(\mathbf{Z}\) -modules \(\mathbf{Z} / p^k\mathbf{Z}\) ( \(k \in \mathbf{N}\) ).
\item
  Let M be the direct sum of the Z-modules \(Z/p^{k}Z\) ( \(k \in N\) ); show that Supp(M) is closed but distinct from V(a), where a is the annihilator of M.
\item
  Let \(\mathbf{N}\) be the direct sum of the \(\mathbf{Z}\) -modules \(\mathbf{Z} / n\mathbf{Z}\) ( \(n \geqslant 1\) ); show that \(\operatorname{Supp}(\mathbf{N})\) is not closed in \(\operatorname{Spec}(\mathbf{Z})\) .
\item
  Deduce from (c) and (d) an example of an A-module M such that, if a is its annihilator, Supp(M) is not closed and its closure is distinct from V(a).
\item
  Show that the support of the \(\mathbf{Z}\) -module \(\prod_{k=1}^{n} (\mathbf{Z}/p^k\mathbf{Z})\) is distinct from the closure of the union of the supports of the factor modules \(\mathbf{Z}/p^k\mathbf{Z}\) .
\end{enumerate}

\begin{enumerate}
\def\labelenumi{\arabic{enumi}.}
\setcounter{enumi}{22}
\item
  Let \(p\) be a prime number; give an example of \(\mathbf{Z}\) -modules \(M, N\) , where \(N\) is finitely generated, such that \(\mathbf{M}_{(p)} \neq 0\) and \(\mathbf{N}_{(p)} \neq 0\) , but \((\mathbf{M} \otimes_{\mathbf{Z}} \mathbf{N})_{(p)} = 0\) (take \(\mathbf{M} = \mathbf{Q}\) ).
\item
  \begin{enumerate}
  \def\labelenumii{(\alph{enumii})}
  \tightlist
  \item
    Let A be a ring and M, N two A-modules such that M is finitely generated. Show that \(\operatorname{Supp}(\operatorname{Hom}_{\mathbf{A}}(\mathbf{M}, \mathbf{N}))\) is contained in the intersection \(\operatorname{Supp}(\mathbf{M}) \cap \operatorname{Supp}(\mathbf{N})\) .
  \end{enumerate}
\end{enumerate}

\begin{enumerate}
\def\labelenumi{(\alph{enumi})}
\setcounter{enumi}{1}
\item
  Give an example where M and N are both finitely presented and \(\operatorname{Supp}(\operatorname{Hom}_{\mathbf{A}}(\mathbf{M},\mathbf{N}))\) is strictly contained in \(\operatorname{Supp}(\mathbf{M}) \cap \operatorname{Supp}(\mathbf{N})\) (take \(\mathbf{A} = \mathbf{Z}_{(p)}\) and \(\mathbf{N} = \mathbf{A}\) ).
\item
  Let M be the Z-module the direct sum of the \(Z/p^{k}Z\) (p prime, \(k \in N\) ). Show that \(\text{Supp}(\text{Hom}_{\mathbf{Z}}(\text{M}, \text{M}))\) is not contained in \(\text{Supp}(\text{M})\) .
\end{enumerate}

\begin{enumerate}
\def\labelenumi{\arabic{enumi}.}
\setcounter{enumi}{24}
\tightlist
\item
  Let A be a Noetherian ring, \(X = \text{Spec}(A)\) its spectrum, M a finitely generated A-module and \(Y = \text{Supp}(M)\) . Let \(Y_k (1 \leqslant k \leqslant n)\) be the distinct connected components of Y.
\end{enumerate}

\begin{enumerate}
\def\labelenumi{(\alph{enumi})}
\item
  Show that there exists a unique decomposition of M as a direct sum \(M = M_{1} \oplus M_{2} \oplus \cdots \oplus M_{n}\) such that \(\text{Supp}(M_{k}) = Y_{k}\) for all k.
\item
  If \(\mathfrak{a} = \operatorname{Ann}(\mathbf{M})\) and \(\mathfrak{a}_k = \operatorname{Ann}(\mathbf{M}_k)\) , show that the \(\mathfrak{a}_k\) are relatively prime in pairs and that \(\prod_{k=1}^{n} \mathfrak{a}_k = \mathfrak{a}\) .
\end{enumerate}

(Using Proposition 17 of no. 4, reduce it to the case where \(\mathfrak{a} = 0\) and \(\mathbf{Y} = \mathbf{X}\) and apply Proposition 15 of no. 3.)

\begin{enumerate}
\def\labelenumi{\arabic{enumi}.}
\setcounter{enumi}{25}
\tightlist
\item
  Let \(\phi: \mathbf{A} \to \mathbf{B}\) be a ring homomorphism such that the mapping
\end{enumerate}

\[
^ {a} \phi : \operatorname{Spec} (\mathbf {B}) \rightarrow \operatorname{Spec} (\mathbf {A})
\]

is surjective. Let N be a finitely generated A-module; show that, if \(u: M \to N\) is an A-module homomorphism such that \(u \otimes 1: M_{(B)} \to N_{(B)}\) is surjective, then u is surjective (apply Proposition 19 of no. 4 to Coker(u)).

\begin{enumerate}
\def\labelenumi{\arabic{enumi}.}
\setcounter{enumi}{26}
\item
  Give an example of a local homomorphism \(\rho: \mathbf{A} \to \mathbf{B}\) and an A-module \(\mathbf{M}\) (not finitely generated) such that \(\operatorname{Supp}(\mathbf{M}_{(\mathbf{B})})\) is not equal to \(a\rho^{-1}(\operatorname{Supp}(\mathbf{M}))\) (take \(\mathbf{A} = \mathbf{Z}_{(p)}\) and \(\mathbf{B} = \mathbf{Z}/p\mathbf{Z}\) for some prime number \(p\) ).
\item
  Give an example of a Z-module M such that, for some prime number p such that \((p) \in \text{Supp}(M)\) , there exists no Z-homomorphism \(\neq 0\) from M to Z/pZ.
\end{enumerate}

